\documentclass[11pt,a4paper]{article}
\usepackage[T1]{fontenc}
\usepackage[utf8]{inputenc}
\usepackage[portuguese]{babel}
\usepackage[margin=2.0cm]{geometry}
\usepackage{enumitem}
\usepackage{titlesec}
\usepackage{xcolor}
\usepackage[hidelinks]{hyperref}
\usepackage{booktabs}
\usepackage{amssymb}

\definecolor{cinza}{gray}{0.35}
\titleformat{\section}{\large\bfseries}{}{0pt}{}[\vspace{-6pt}\rule{\textwidth}{0.4pt}]
\titlespacing{\section}{0pt}{9pt}{4pt}
\newcommand{\lin}[2]{\noindent\textbf{#1}\hfill{\color{cinza}\small #2}\par}
\setlist[itemize]{leftmargin=1.2em,itemsep=1pt,topsep=2pt,parsep=0pt}
\pagestyle{empty}

\begin{document}

% CC BY 4.0 · Copyright (c) 2026 Aarão Melo Lopes · ver LICENSE na raiz.

\begin{center}
{\LARGE\bfseries Aarão Melo Lopes}\\[2pt]
{\color{cinza}\emph{Engenheiro Eletricista \,|\, Matemática Aplicada \& Engenharia de Sistemas}}\\[2pt]
{\small
\href{mailto:aarao.melo.lopes@gmail.com}{aarao.melo.lopes@gmail.com}
\quad·\quad
\href{https://goldenkingdom.patriatechnology.com}{goldenkingdom.patriatechnology.com}}
\end{center}

\vspace{-2pt}
\section*{O que faço e o que procuro}

Trabalho na fronteira entre álgebra e sistemas: a dualidade como memória da divisão ---
guardar a metade que a divisão deitaria fora. Em três volumes: a \textbf{Teoria} afirma;
o \textbf{Catálogo} mede corpo a corpo (\emph{o que não é citado não é testado}); o
\textbf{Enredo} conta a mesma estrutura sem uma fórmula. Resultado actual: 61 + 457 + 365\,pp.,
328 testes a zero, seis sistemas. Procuro parceiros onde ainda não corre sozinho.

\section*{Em que pé está cada coisa}

\begin{center}\footnotesize
\begin{tabular}{@{}p{2.3cm}p{8.0cm}p{4.2cm}@{}}
\toprule
\textbf{estado} & \textbf{o que é} & \textbf{onde} \\
\midrule
\textbf{demonstrado} &
21 enunciados com prova escrita &
\emph{Teoria}, 58\,pp. \\[1pt]
\textbf{medido} &
328 testes em C; a bateria sai a zero &
\emph{Catálogo} e \texttt{tests/} \\[1pt]
\textbf{implementado} &
ISA de uma instrução (\texttt{MOVE}); banco sem \emph{malloc};
assistente sem modelo externo; compilador e motor de xadrez &
\texttt{banco/}, \texttt{broca-so}, \texttt{chess} \\[1pt]
\textbf{proposto} &
liga grafeno--estanho; Headjack; CristalChain; jogo a partir do GDD &
propostas abaixo \\
\bottomrule
\end{tabular}
\end{center}

\section*{Produção}

\lin{Três volumes --- teoria, medição, narrativa}{2026, em curso}
\begin{itemize}
\item \textbf{Teoria} (61\,pp.) --- 21 enunciados e 21 provas; uma definição (corpo dual).
\item \textbf{Catálogo} (457\,pp.) --- bestiário: cada corpo com a sua régua e a sua dinâmica;
medidores citados nominalmente. É o \emph{embaixo}: onde afirmação e narrativa se encontram
e se medem.
\item \textbf{Enredo} (365\,pp.) --- a mesma estrutura como mito de xadrez: Cristal (23 lances)
e Quasicristal (128 cantos do Cruzeiro); semente no \emph{gambito recusado}.
\item \textbf{328 medidores em C} --- 328 verdes, 0 falhas; 131 só com aritmética inteira.
\textbf{O método:} não se compara contra um valor esperado --- isso é ler metade, e metade dissipa
--- \emph{reverte-se e lê-se o resíduo}, resolvendo e tirando a prova na mesma passagem. E a medida
são as \textbf{duas metades}: o resíduo onde tem de ser zero diz que \emph{existe}, o resíduo onde
não pode ser zero diz que é \emph{único}. Cada asserção é ainda submetida a \emph{mutação} --- se
estragar o código não a faz falhar, ela não media nada.
\item \textbf{Dois artigos auto-contidos} --- \emph{O corpo estelar} (19\,pp.), a especificação do
sistema: de uma única exigência derivam-se o coeficiente da equação de campo, a dimensão quatro, o
termo cosmológico, o factor $8\pi$ e as constantes em função da régua. E \emph{Dual Sort} (6\,pp.),
a ordenação sem comparações --- onde se mostra que \textbf{ordenar e cortar são duais} e que o
\emph{corte máximo} (\emph{max-cut}) numa árvore sai por leitura, não por busca.
\end{itemize}

\smallskip
\lin{Sistemas construídos}{abr.--ago.\ 2026}
Compilador, VM de pilha, núcleo, orquestração de agentes (\texttt{broca-so}), motor de
xadrez e corpo teórico. Publicação com portão de verificação.
\\[1pt]
{\small\color{cinza}C · Makefile · Shell · POSIX · Git · \texttt{MOVE} · I/O sem \emph{malloc}
· assembly/\texttt{\_start}}

\section*{Formação}

\lin{Mestrado em Ciência da Computação --- IME-USP}{2020 -- 2022}
Orientador: Prof.\ Dr.\ Ronaldo Fumio Hashimoto. Bolsista CAPES.
\textbf{72 créditos} (exigência: 44); conceito máximo em oito das nove disciplinas.
O corpo teórico e os sistemas substituíram a dissertação formal --- disponível sob consulta.

\smallskip
\lin{Extensão --- IME-USP}{verão 2021}
Cálculo no $\mathbb{R}^n$ e Espaços Métricos (60\,h cada); frequência 100\,\%.

\smallskip
\lin{Bacharelado em Engenharia Elétrica --- UFRR}{2011 -- 2018}
Média \textbf{8{,}66} · 4\,201 horas. Nota 10 em Cálculo, Eletromagnetismo, Conversão de Energia;
acima de 9 em Álgebra Linear, Cálculo Numérico, Estatística, Métodos Matemáticos.

\smallskip
\lin{\emph{TCC}: Controle Direto de Torque do Motor de Indução Trifásico}{UFRR, 2018}
123\,pp. Nota 10 da banca, recomendação de publicação. A tabela de comutação é, em 2026,
o corpo mínimo da teoria (\texttt{MOVE}, sem estado).

\smallskip
\lin{Monitoria --- 473 horas}{UFRR, 2013 e 2016--17}
Cálculo I (108\,h) e Álgebra Linear I no curso de Matemática (365\,h).

\section*{Propostas --- onde procuro parceiros}

\lin{\textbf{I.} Headjack + liga grafeno--estanho}{engenharia}
Formalismo e máquina já correm. Liga em \texttt{tests/liga.c}, Headjack em
\texttt{tests/headjack.c}: campo, núcleo do operador, limites do que se vê.
\\[1pt]
\textbf{O que preciso do parceiro:} bancada para a liga; engenharia para agentes em produção.

\smallskip
\lin{\textbf{II.} Mundo do Enredo --- jogo e cinema}{criação}
365\,pp., 26 partes, 148 capítulos; GDD no catálogo. Vitória: o \emph{mate que
não mata} --- reacoplar as duas mãos, resíduo zero, relógio a correr.
\\[1pt]
\textbf{O que preciso do parceiro:} equipa de jogo/cinema; trago o mundo e a coerência.

\smallskip
\lin{\textbf{III.} CristalChain --- livro-razão cuja cifra conserva}{distribuído}
Cifra forçada: ouro branco ($\det=+1$) conserva; a família metálica ($\det=-1$) anti-conserva.
\textbf{O que esta proposta abre:} a álgebra e os medidores. \textbf{O que ela não inclui:}
rede, nós e consenso --- \emph{o que se propõe aqui é a cifra e a paragem, não um sistema
distribuído}, e por si só isto não é uma blockchain nem resolve dupla despesa.
\\[1pt]
\textbf{O que preciso do parceiro:} sistemas distribuídos / consenso.

\section*{Colaboração}

Meço em vez de argumentar. \textbf{Presencial:} grupo, bancada e prazo --- escreva-me.
\textbf{Remota:} CC BY 4.0 / MIT; pode avançar sem combinar.

\medskip
\begin{center}\small
\textbf{Documentos e código}\\[3pt]
\href{https://goldenkingdom.patriatechnology.com/docs/teoria.pdf}{teoria.pdf}
\,·\,
\href{https://goldenkingdom.patriatechnology.com/docs/catalogo.pdf}{catalogo.pdf}
\,·\,
\href{https://goldenkingdom.patriatechnology.com/docs/enredo.pdf}{enredo.pdf}
\,·\,
\href{https://goldenkingdom.patriatechnology.com/docs/livro.pdf}{livro.pdf}
\\[2pt]
\href{https://goldenkingdom.patriatechnology.com/repo.git}{\texttt{git clone https://goldenkingdom.patriatechnology.com/repo.git}}
\\[3pt]
{\footnotesize\color{cinza}CC BY 4.0 (textos) · MIT (código)}
\end{center}

\end{document}
